% !TEX root = ../main.tex
% Generated from full-precision results/summary/f1_raw_scale; display only.
% F1 numbers: percent; differences: percentage points. Units are in the prose.
\newcommand{\FWithinCornMin}{90.98}
\newcommand{\FWithinCornMax}{96.23}
\newcommand{\FCrossCornMin}{85.68}
\newcommand{\FCrossCornMax}{92.92}
\newcommand{\FWithinSoyMin}{89.55}
\newcommand{\FWithinSoyMax}{95.89}
\newcommand{\FCrossSoyMin}{83.08}
\newcommand{\FCrossSoyMax}{92.15}
\newcommand{\FModelGapCMin}{1.43}
\newcommand{\FModelGapCMax}{1.61}
\newcommand{\FModelGapDMin}{3.34}
\newcommand{\FModelGapDMax}{3.98}
\newcommand{\FCUtaeUnetCorn}{0.04}
\newcommand{\FCUtaeUnetSoy}{0.08}
\newcommand{\FConvCrossCorn}{88.73}
\newcommand{\FConvGapCorn}{2.79}
\newcommand{\FConvCrossSoy}{87.64}
\newcommand{\FConvGapSoy}{3.00}
\newcommand{\FOtherCrossMin}{89.69}
\newcommand{\FOtherCrossMax}{90.39}
\newcommand{\FOtherGapMin}{3.74}
\newcommand{\FOtherGapMax}{3.97}
\newcommand{\FCrossModelGapMin}{1.45}
\newcommand{\FCrossModelGapMax}{2.13}
\newcommand{\FConvSoyBD}{7.39}
\newcommand{\FConvSoyDB}{3.36}
\newcommand{\FConvSoyCD}{0.90}
\newcommand{\FUnetSoyBD}{4.02}
\newcommand{\FUnetSoyDB}{6.14}
\newcommand{\FUnetSoyCD}{3.54}
\newcommand{\FConvSoyCDRunMin}{0.77}
\newcommand{\FConvSoyCDRunMax}{0.99}
\newcommand{\FConvSoyBDRunMin}{3.29}
\newcommand{\FConvSoyBDRunMax}{10.81}

% Derived full-precision values; selectors and formulas in sec31_joint_boundary/prose_evidence.csv.
\newcommand{\SACropCorn}{22.42}
\newcommand{\SABoundCorn}{16.90}
\newcommand{\SACropSoy}{19.45}
\newcommand{\SABoundSoy}{17.42}
\newcommand{\SAInteriorMax}{0.64}
\newcommand{\SBSmallArea}{23.76}
\newcommand{\SBSmallFnMin}{36.89}
\newcommand{\SBSmallFnMax}{37.90}
\newcommand{\SBSmallBound}{35.42}
\newcommand{\SBSmallBoundaryFn}{43.21}
\newcommand{\SBSmallInteriorFn}{3.18}
\newcommand{\SBComplexArea}{0.49}
\newcommand{\SBComplexBound}{69.29}
\newcommand{\SDSmallCorn}{12.91}
\newcommand{\SDSmallSoy}{14.97}
\newcommand{\SDMainArea}{11.19}
\newcommand{\SDReductionBoundary}{12.45}
\newcommand{\SDReductionInterior}{2.04}
\newcommand{\SDNetShare}{20.85}
\newcommand{\SCMainArea}{8.15}
\newcommand{\SCReductionBoundary}{1.95}
\newcommand{\SCReductionInterior}{1.35}
\newcommand{\SCNetShare}{10.22}
\newcommand{\SDSoyNetCorn}{19.67}
\newcommand{\SDSoyNetSoy}{17.68}
\newcommand{\SCUtaeBoundary}{25.79}
\newcommand{\SCUnetBoundary}{28.03}
\newcommand{\SCUtaeInterior}{0.93}
\newcommand{\SCUnetInterior}{0.81}
\newcommand{\SCNegativeArea}{8.63}
\newcommand{\SCNegativePp}{0.18}


\section{实验结果与分析}

\subsection{模型站内分类性能}\label{sec:within_site}

% 数据：results/summary/f1_raw_scale/within_site_crop_f1.csv（汇总口径）、within_site_crop_f1_fold_stats.csv（折间标准差）
% F1以百分数、差值以百分点显示两位小数；正文数值由generate_core_tables.py从完整精度生成。

如图\ref{fig:within_site_f1}所示，三种模型在四个站点的站内玉米F1为\FWithinCornMin{}\%～\FWithinCornMax{}\%，大豆F1为\FWithinSoyMin{}\%～\FWithinSoyMax{}\%。三种模型的分类表现呈现出两点共性：在四个站点上，玉米F1均高于大豆F1，差值均小于2.00个百分点；跨站点比较，各模型的玉米和大豆F1均在A站点取得各自最高值。五折交叉验证中，同一数据集不同划分下的每折独立训练，12组实验（3种模型×4个站点）的玉米和大豆F1折间标准差均小于1.00\%。

\par\addvspace{3pt}
\noindent\begin{minipage}{\linewidth}
\captionsetup{type=figure}
\centering
\begin{subfigure}{\linewidth}
\centering
\includegraphics{results/fig_within_site_f1_a.pdf}
\caption{四站点玉米分类精度}
\end{subfigure}\par\nointerlineskip\vspace{3pt}
\begin{subfigure}{\linewidth}
\centering
\includegraphics{results/fig_within_site_f1_b.pdf}
\caption{四站点大豆分类精度}
\end{subfigure}\par\nointerlineskip\vspace{3pt}
\modellegend
\caption{三种模型在四个站点的站内玉米和大豆分类精度}
\label{fig:within_site_f1}
\end{minipage}\par\addvspace{3pt}

三种模型之间的差异主要体现在ConvGRU与另两种模型之间。U-TAE和3D U-Net在各站点的玉米和大豆F1均高于ConvGRU，而二者之间差值均小于1.00个百分点，互有高低：U-TAE在A、B两个站点略优，3D U-Net在D站点略优，C站点U-TAE的玉米和大豆F1分别高出\FCUtaeUnetCorn{}和\FCUtaeUnetSoy{}个百分点。ConvGRU与前两者的差距因站点而异：两种作物在C站点为\FModelGapCMin{}～\FModelGapCMax{}个百分点，在D站点为\FModelGapDMin{}～\FModelGapDMax{}个百分点。

结合1.3节的面积—形状分布，进一步按完整参考斑块分组考察漏分位置，图\ref{fig:within_site_errors}给出三模型的全站作物对错分布。A站两种作物面积占比最大的组合均为100～200~ha、形状指数1.5～2，分别占全站玉米和大豆面积的\SACropCorn{}\%和\SACropSoy{}\%。这两个组合的边界占比均低于对应作物的全站水平，三模型的内部漏分率最高为\SAInteriorMax{}\%，边界漏分率也低于同站对应作物的全站边界水平。主要面积所在组合的边界和内部均有较好的识别表现，与A站两种作物F1均最高相对应。

B站大豆较小斑块中的漏分，主要来自有较大面积支持的较规则形状组。面积小于50~ha、形状指数小于1.5的斑块占该站大豆面积的\SBSmallArea{}\%，包含三模型全站大豆漏分的\SBSmallFnMin{}\%～\SBSmallFnMax{}\%。该组边界占比为\SBSmallBound{}\%，ConvGRU在边界和内部的漏分率分别为\SBSmallBoundaryFn{}\%和\SBSmallInteriorFn{}\%，说明漏分量既与边界所占面积有关，也与边界内的识别不足对应。形状指数不小于3的小斑块虽有更高的边界占比和组内漏分率，却仅占大豆面积的\SBComplexArea{}\%，不是全站漏分的主要来源。D站两种作物在小于50~ha、形状指数小于1.5组中的面积占比分别为\SDSmallCorn{}\%和\SDSmallSoy{}\%，组内边界占比和内部漏分率也接近，与两作物分类表现接近相对应，但不能仅据面积或形状分布判断识别难度。

模型间漏分差异在C、D站对应的组合也不同。D站玉米10～50~ha、形状指数小于1.5组占玉米面积的\SDMainArea{}\%，3D U-Net相对ConvGRU在该组边界和内部的漏分率分别降低\SDReductionBoundary{}和\SDReductionInterior{}个百分点，二者合计占全站玉米漏分净减少量的\SDNetShare{}\%；该组也是D站大豆两种比较模型漏分净减少量最大的组合。C站玉米中，3D U-Net漏分净减少量最大的组合为200～500~ha、形状指数不小于3，占玉米面积的\SCMainArea{}\%，边界和内部漏分率分别降低\SCReductionBoundary{}和\SCReductionInterior{}个百分点，占全站玉米漏分净减少量的\SCNetShare{}\%。D站这一主要组合的改善幅度和贡献均大于C站上述组合，但两站其余组合仍共同影响全站差距。ConvGRU的内部漏分也并非只集中于复杂斑块：A站两种作物、B站大豆和D站两种作物的内部漏分计数最多的组合均为10～50~ha、形状指数小于1.5，而C站玉米为上述大型复杂组合。

U-TAE与3D U-Net在边界和内部的优势并不总是一致。C站玉米的全站边界漏分率，U-TAE为\SCUtaeBoundary{}\%，低于3D U-Net的\SCUnetBoundary{}\%；内部则为\SCUtaeInterior{}\%，高于3D U-Net的\SCUnetInterior{}\%，D站大豆也存在同样的位置差异。模型改善还存在反例：C站面积不小于500~ha、形状指数不小于3的玉米斑块占该站玉米面积的\SCNegativeArea{}\%，U-TAE在其内部的漏分率比ConvGRU高\SCNegativePp{}个百分点。图\ref{fig:within_site_errors}用于观察全站错分的空间聚集，细小边界及组内差异由上述全站统计定位。这些面积—形状与错误位置的对应不表示独立因果作用；作物漏分分析也未包含参考其它地类误分入作物的影响，不能直接分解F1差异。

\begin{figure*}[!t]
\centering
% Native 5.2 cm panels; full 2560x2560 reference extent. No rescaling.
\begin{subfigure}{\dimexpr(\textwidth-6pt)/3\relax}\centering
\includegraphics{results/fig_within_site_full_errors_A_convgru.pdf}
\caption{A站ConvGRU作物对错图}
\end{subfigure}\hspace{3pt}%
\begin{subfigure}{\dimexpr(\textwidth-6pt)/3\relax}\centering
\includegraphics{results/fig_within_site_full_errors_A_unet3d.pdf}
\caption{A站3D U-Net作物对错图}
\end{subfigure}\hspace{3pt}%
\begin{subfigure}{\dimexpr(\textwidth-6pt)/3\relax}\centering
\includegraphics{results/fig_within_site_full_errors_A_utae.pdf}
\caption{A站U-TAE作物对错图}
\end{subfigure}\par\nointerlineskip\vspace{3pt}
\begin{subfigure}{\dimexpr(\textwidth-6pt)/3\relax}\centering
\includegraphics{results/fig_within_site_full_errors_B_convgru.pdf}
\caption{B站ConvGRU作物对错图}
\end{subfigure}\hspace{3pt}%
\begin{subfigure}{\dimexpr(\textwidth-6pt)/3\relax}\centering
\includegraphics{results/fig_within_site_full_errors_B_unet3d.pdf}
\caption{B站3D U-Net作物对错图}
\end{subfigure}\hspace{3pt}%
\begin{subfigure}{\dimexpr(\textwidth-6pt)/3\relax}\centering
\includegraphics{results/fig_within_site_full_errors_B_utae.pdf}
\caption{B站U-TAE作物对错图}
\end{subfigure}\par\nointerlineskip\vspace{3pt}
\begin{subfigure}{\dimexpr(\textwidth-6pt)/3\relax}\centering
\includegraphics{results/fig_within_site_full_errors_C_convgru.pdf}
\caption{C站ConvGRU作物对错图}
\end{subfigure}\hspace{3pt}%
\begin{subfigure}{\dimexpr(\textwidth-6pt)/3\relax}\centering
\includegraphics{results/fig_within_site_full_errors_C_unet3d.pdf}
\caption{C站3D U-Net作物对错图}
\end{subfigure}\hspace{3pt}%
\begin{subfigure}{\dimexpr(\textwidth-6pt)/3\relax}\centering
\includegraphics{results/fig_within_site_full_errors_C_utae.pdf}
\caption{C站U-TAE作物对错图}
\end{subfigure}\par\nointerlineskip\vspace{3pt}
\begin{subfigure}{\dimexpr(\textwidth-6pt)/3\relax}\centering
\includegraphics{results/fig_within_site_full_errors_D_convgru.pdf}
\caption{D站ConvGRU作物对错图}
\end{subfigure}\hspace{3pt}%
\begin{subfigure}{\dimexpr(\textwidth-6pt)/3\relax}\centering
\includegraphics{results/fig_within_site_full_errors_D_unet3d.pdf}
\caption{D站3D U-Net作物对错图}
\end{subfigure}\hspace{3pt}%
\begin{subfigure}{\dimexpr(\textwidth-6pt)/3\relax}\centering
\includegraphics{results/fig_within_site_full_errors_D_utae.pdf}
\caption{D站U-TAE作物对错图}
\end{subfigure}\par\nointerlineskip\vspace{3pt}
\definecolor{sec31FullCorrect}{HTML}{377EB8}
\definecolor{sec31FullError}{HTML}{D73027}
{\songti\zihao{-5}\legbox{sec31FullCorrect}\,正确作物像元\quad\legbox{sec31FullError}\,错分作物像元\quad\legbox{cOther}\,参考其它地类\qquad
\begin{tikzpicture}[x=0.677083333cm,y=1cm,baseline=-0.7ex]
\draw[line width=0.6pt] (0,0)--(2,0);
\foreach \x in {0,1,2} {\draw[line width=0.6pt] (\x,0)--(\x,0.07);}
\node[anchor=north,font=\songti\zihao{-5},inner sep=1pt] at (0,0) {0};
\node[anchor=north,font=\songti\zihao{-5},inner sep=1pt] at (1,0) {10};
\node[anchor=north west,font=\songti\zihao{-5},inner sep=1pt] at (2,0) {20 km};
\end{tikzpicture}\par}

\caption{四站三种模型的全站作物参考像元错分分布}
\label{fig:within_site_errors}
\par\noindent\begin{minipage}{\textwidth}\songti\zihao{-5}
注：行对应A—D站，列依次为ConvGRU、3D U-Net和U-TAE；同站三图覆盖范围一致。红色包括参考玉米、大豆被判为另一作物或其它地类；参考其它地类始终为灰色，不显示其中误分入作物的像元。各图保留原始像元，细小边界差异以全站分组统计为依据。
\end{minipage}
\end{figure*}

\subsection{模型跨站分类性能与泛化退化}

% 数据：tables/table2_cross_site_f1_12paths.csv ← results/summary/f1_raw_scale/cross_site_crop_f1_summary.csv
%       results/summary/f1_raw_scale/within_site_crop_f1.csv（同模型目标站站内基准）
%       results/summary/f1_raw_scale/crop_f1_difference_36paths.csv（图 4，原尺度差值）

如表\ref{tab:cross_site_f1_paths}所示，源站模型不经微调直接用于目标站后，36个模型—路径组合的玉米和大豆跨站F1均低于同一模型在目标站的站内汇总F1。跨站玉米F1为\FCrossCornMin{}\%～\FCrossCornMax{}\%，大豆F1为\FCrossSoyMin{}\%～\FCrossSoyMax{}\%，每项为五次完整目标站预测的均值。按目标站对3条路径取平均，三种模型的玉米和大豆跨站F1均在以A站点为目标站时最高，与站内结果中A站点精度最高相对应；玉米跨站F1则均在以C站点为目标站时最低，大豆跨站F1的最低目标站则不一致。

\begin{table*}[tbp]
\centering
\caption{三种模型在12条迁移路径上的跨站玉米和大豆F1}
\label{tab:cross_site_f1_paths}
{\songti\zihao{-5}
\renewcommand{\arraystretch}{1.12}
\setlength{\tabcolsep}{7pt}
\begin{tabular*}{\textwidth}{@{\extracolsep{\fill}}c*{6}{c}}
\toprule
\multirow{2}{*}{迁移路径} & \multicolumn{2}{c}{ConvGRU} & \multicolumn{2}{c}{3D U-Net} & \multicolumn{2}{c}{U-TAE} \\
\cmidrule(lr){2-3}\cmidrule(lr){4-5}\cmidrule(lr){6-7}
 & 玉米F1（\%） & 大豆F1（\%） & 玉米F1（\%） & 大豆F1（\%） & 玉米F1（\%） & 大豆F1（\%） \\
\midrule
A→B & 89.91 ± 0.11 & 87.52 ± 0.40 & 91.43 ± 0.39 & 89.28 ± 0.40 & 92.09 ± 0.26 & 89.90 ± 0.38 \\
A→C & 86.97 ± 0.59 & 86.94 ± 1.20 & 87.71 ± 0.59 & 87.96 ± 0.34 & 88.20 ± 0.61 & 88.29 ± 0.69 \\
A→D & 88.54 ± 0.67 & 88.05 ± 1.16 & 90.64 ± 0.53 & 90.56 ± 0.45 & 91.78 ± 0.43 & 91.48 ± 0.41 \\
B→A & 92.00 ± 0.18 & 91.59 ± 0.13 & 92.92 ± 0.59 & 92.15 ± 0.44 & 92.70 ± 0.25 & 91.87 ± 0.31 \\
B→C & 85.68 ± 0.95 & 84.52 ± 1.63 & 88.09 ± 0.55 & 88.20 ± 0.81 & 87.70 ± 0.58 & 87.66 ± 0.84 \\
B→D & 85.73 ± 1.67 & 83.08 ± 2.70 & 90.87 ± 0.33 & 90.43 ± 0.41 & 89.59 ± 1.41 & 88.50 ± 2.20 \\
C→A & 90.31 ± 0.59 & 90.15 ± 0.49 & 92.09 ± 0.40 & 91.64 ± 0.43 & 91.52 ± 0.61 & 91.09 ± 0.39 \\
C→B & 87.47 ± 1.05 & 85.06 ± 0.96 & 90.23 ± 1.00 & 88.98 ± 0.88 & 90.33 ± 0.62 & 88.18 ± 0.52 \\
C→D & 89.77 ± 0.29 & 89.57 ± 0.08 & 90.77 ± 0.85 & 90.91 ± 0.17 & 90.98 ± 0.39 & 90.70 ± 0.48 \\
D→A & 91.05 ± 0.24 & 90.85 ± 0.52 & 91.36 ± 0.68 & 91.32 ± 0.68 & 91.45 ± 0.26 & 92.04 ± 0.37 \\
D→B & 89.44 ± 0.24 & 86.19 ± 0.46 & 88.13 ± 1.24 & 86.19 ± 1.91 & 90.45 ± 0.45 & 89.17 ± 0.98 \\
D→C & 87.85 ± 0.30 & 88.19 ± 0.54 & 87.88 ± 0.71 & 88.60 ± 0.83 & 87.83 ± 0.89 & 88.41 ± 0.84 \\
\midrule
平均 & 88.73 & 87.64 & 90.18 & 89.69 & 90.39 & 89.77 \\
\bottomrule

\end{tabular*}
\par\nointerlineskip\vspace{3pt}
\begin{minipage}{\textwidth}
\noindent 注：箭头左侧为源站，右侧为目标站。每格为源站5个折模型分别预测目标站全部400个图块所得5次F1的算术平均，以百分数表示，未进行微调和模型集成；“平均”行为12条迁移路径F1均值的算术平均。
\end{minipage}
}
\end{table*}

% 第二段：12 条路径平均跨站 F1 与平均 F1 差值均从完整精度计算。
% 两种排序分别对应实际精度与相对同一模型目标站当地训练基准的差距。

按12条路径取算术平均，ConvGRU的跨站F1均值最低，但平均F1差值最小：其跨站玉米和大豆F1分别为\FConvCrossCorn{}\%和\FConvCrossSoy{}\%，比3D U-Net和U-TAE低\FCrossModelGapMin{}～\FCrossModelGapMax{}个百分点（后两者两种作物的均值为\FOtherCrossMin{}\%～\FOtherCrossMax{}\%）；而玉米和大豆平均F1差值分别为\FConvGapCorn{}和\FConvGapSoy{}个百分点，小于3D U-Net和U-TAE（\FOtherGapMin{}～\FOtherGapMax{}个百分点）。跨站F1比较的是模型在目标站的实际分类精度，F1差值衡量的则是跨站F1与同一模型在目标站站内汇总F1之间的差距。因此，F1差值较小不代表跨站绝对精度更高，评价模型的跨站表现需要同时考察这两个指标。

% 第三段：图 4 数据 results/summary/f1_raw_scale/crop_f1_difference_36paths.csv。
% 固定目标站D的单次差值由未舍入本站基准减每次跨站F1生成，不由显示值反算。
% B↔D：ConvGRU 玉米、大豆均 B→D 较大，3D U-Net 均 D→B 较大；比较基于各路径五次预测均值。

如图\ref{fig:degradation_heatmap}所示，F1差值随源—目标组合而变化，且三种模型的变化并不一致。即使目标站相同，F1差值也会随源站不同而改变：以D站点为目标站时，ConvGRU的大豆F1差值在源站为C时为\FConvSoyCD{}个百分点，源站为B时为\FConvSoyBD{}个百分点；这两种源站下五次预测对应的F1差值分别为\FConvSoyCDRunMin{}～\FConvSoyCDRunMax{}和\FConvSoyBDRunMin{}～\FConvSoyBDRunMax{}个百分点，源站为B时的每一次结果均高于源站为C时，但五次之间的波动也更大。在B、D两个站点之间，两个方向上的F1差值大小关系也因模型而异：ConvGRU的玉米和大豆F1差值均为B→D方向较大，3D U-Net则均为D→B方向较大；以大豆为例，ConvGRU在B→D和D→B路径上分别为\FConvSoyBD{}和\FConvSoyDB{}个百分点，3D U-Net分别为\FUnetSoyBD{}和\FUnetSoyDB{}个百分点。上述方向差异基于五次预测的均值，仅反映这一站点对上的表现，不代表普遍的迁移规律。这些差异与站点空间特征的对应关系，将在后文结合地物组成与景观结构分析。

\begin{figure*}[tbp]
\centering
\begin{subfigure}{\dimexpr(\textwidth-6pt)/3\relax}\centering
\includegraphics{results/fig_degradation_heatmap_a.pdf}
\caption{ConvGRU玉米F1差值}
\end{subfigure}\hspace{3pt}%
\begin{subfigure}{\dimexpr(\textwidth-6pt)/3\relax}\centering
\includegraphics{results/fig_degradation_heatmap_b.pdf}
\caption{3D U-Net玉米F1差值}
\end{subfigure}\hspace{3pt}%
\begin{subfigure}{\dimexpr(\textwidth-6pt)/3\relax}\centering
\includegraphics{results/fig_degradation_heatmap_c.pdf}
\caption{U-TAE玉米F1差值}
\end{subfigure}\par\nointerlineskip\vspace{3pt}
\begin{subfigure}{\dimexpr(\textwidth-6pt)/3\relax}\centering
\includegraphics{results/fig_degradation_heatmap_d.pdf}
\caption{ConvGRU大豆F1差值}
\end{subfigure}\hspace{3pt}%
\begin{subfigure}{\dimexpr(\textwidth-6pt)/3\relax}\centering
\includegraphics{results/fig_degradation_heatmap_e.pdf}
\caption{3D U-Net大豆F1差值}
\end{subfigure}\hspace{3pt}%
\begin{subfigure}{\dimexpr(\textwidth-6pt)/3\relax}\centering
\includegraphics{results/fig_degradation_heatmap_f.pdf}
\caption{U-TAE大豆F1差值}
\end{subfigure}\par\nointerlineskip\vspace{3pt}
\includegraphics{results/fig_degradation_heatmap_cbar.pdf}\par\nointerlineskip\vspace{3pt}
{\songti\zihao{-5}F1下降（百分点）\par}
\caption{不同源—目标区域组合下三种模型的玉米和大豆F1差值（单位：百分点；行为源站，列为目标站，对角线不适用；同模型目标站站内F1减跨站五次均值，六幅子图共用色标）}
\label{fig:degradation_heatmap}
\end{figure*}

\FloatBarrier
\subsection{模型跨站泛化差异的空间解释}\label{sec:spatial_explain}
